\histitem{局部紧群的表示}

直到1939年，拓扑群表示的研究局限于紧情形和交换情形。另一方面，物理学家在量子力学中发现了许多理由，需要研究不可约酉表示，以寻找能够模拟基本粒子状态的希尔伯特空间。P. 狄拉克在 E. 马约拉纳工作的基础上，于1936年前后指出，在这一视角下，洛伦兹群 SO(3,1) 和庞加莱群 SO(3,1) × \(R^{4}\) 的意义。他们的方法仍与当时数学家的方法相去甚远。但1939年，E. 维格纳 {[}90{]} 以默里和冯·诺伊曼的工作为基础，对庞加莱群的不可约酉表示作出分类，打开了一道缺口。

我们似乎应将研究局部紧群表示的一般路径归功于盖尔范德。他在1942年写的一篇回忆录中（参见 {[}23{]}，第II卷，第3--17页）与拉伊科夫注意到，这一理论之所以可能，是因为酉表示与正定函数之间存在联系。他们证明每个正定函数都有希尔伯特实现（V，第432页定理1），并由此推出不可约酉表示分离 G 的点（V，第454页定理4）。对合代数 L \(^{1}\) (G) 上的正线性泛函发挥了基本作用：盖尔范德与拉伊科夫揭示了它们同正定函数相联系的关系（参见 V，第448页命题13）。对合代数的包络 C* 代数在 {[}24，第48节{]} 中已经（未命名）出现，不过当时还没有直接提出把代数 Stell(G) 与 G 联系起来（参见 I 第125页定义9）。

1945年以后，酉表示的研究突飞猛进。与 C*-代数的互动以及对量子理论基础的分析很快带来有价值的一般结果：约1947年，西格尔把各种 \(^{(9)}\) 星代数同一个局部紧群联系起来 {[}67{]}，并编织起它们的表示与所研究群的表示之间的联系。不久后，G. 麦基提炼出诱导表示的概念 {[}42{]}，它将在酉表示的具体结果中无处不在。

与此同时，对揭示其多样性与深度的例子的研究大大推动了该理论。盖尔范德和奈马克研究 SL(2, C) 这一例子，于1947年分离出诱导表示概念的一个特殊情形，随后证明它是研究该群表示的关键 {[}23，第II卷，第41--124页{]}。同年，V. 巴格曼研究群 SL(2, R)，发现一个可数的平方可积表示族，称为“离散系列”，以及它们的矩阵系数之间的正交关系（参见 V，第424页命题8） {[}5{]}。正如戈德芒立即认识到的 {[}25{]}，Bargmann 正交关系由所有局部紧群的平方可积表示满足。

很快，对特殊群类的研究又一次、并且长时期内优先于抽象理论。特别是，哈里什-钱德拉将着手一般研究约化李群的表示——这项史诗般的任务，以及朗兰兹将为数论预见的卓越后果，我们都无法在此描述。

